\section{Three errors that must not be confused}
\label{l2:errors}
\noindent\textbf{Aim.} Turn a finite element Hessian entry into an enclosure
of the continuous entry. Then investigate the transmission regularity
behind its observed quadratic convergence. A 90-minute lecture can cover
the error identities and first-derivative regularity; the second-derivative
argument below is suitable for a further reading session.

The geometry is the same regular polygon and coarse fan as in the first
lecture. Write $a(v,w)=\int_P\nabla v\cdot\nabla w$ and
$\energy v=\|\nabla v\|_{L^2(P)}$. A displacement $q$ has the fan extension
$\theta_q=\sum_jq_j\phi_j$. These are vertex displacements, not functions
that can be chosen independently on adjacent triangles.
Recall that $-\Delta u=1$, $u\in H^1_0(P)$, $\ell(v)=\int_Pv$,
$J=\frac12\int_Pu$, and $\F=J/A^2$ with $A=|P|$.

There are three distinct objects:
\begin{center}
\begin{tabular}{lll}\toprule
object & meaning & source of the difference\\\midrule
$u$ & continuous weak solution & ---\\
$u_h$ & exact Galerkin solution in $V_h$ & spatial discretization\\
$\widetilde u$ & computed coefficient vector in $V_h$ & linear-system solve\\\bottomrule
\end{tabular}
\end{center}
Interval rounding is a third error, incurred while evaluating quantities
in the proof. It is handled by enclosing each arithmetic operation.
A small solver residual alone controls neither $u-u_h$ nor the Hessian
error. Conversely, an a priori interpolation theorem does not control the
rounding error in a computed residual.

Each coarse triangle is split uniformly into $m^2$ fine triangles. The
mesh has $nm^2$ triangles and $1+nm(m+1)/2$ vertices. Its boundary has
$nm$ vertices. For fixed $n$, $h=m^{-1}$ is proportional to the maximum
fine-element diameter. All rays are unions of element edges, and the
coefficients obtained by differentiating the pullback are constant on each
fine triangle. A mesh crossing a ray cannot use the same broken
interpolation argument without additional analysis.
The space $V_h$ consists of continuous functions that are affine on each
fine triangle and zero on the polygon boundary. The exact Galerkin state
is characterized by $a(u_h,v_h)=\ell(v_h)$ for every $v_h\in V_h$.
It is a mathematical solution of a finite linear system, independently
of how accurately a numerical solver finds its coefficients.

\section{A residual is an error measured in the dual norm}
\label{l2:residual}
For $u$ satisfying $a(u,v)=\ell(v)$ and any conforming candidate $w$,
\[
 a(u-w,v)=\ell(v)-a(w,v).
\]
Taking the supremum over $\energy v=1$ gives an equality:
\[
 \energy{u-w}=\|\ell-a(w,\cdot)\|_{(H^1_0,a)^*}.
\]
The equality does not compute the dual norm. It identifies what a guaranteed
majorant must bound.

For torsion, let $y\in H(\operatorname{div};P)$ satisfy
$\operatorname{div}y=-1$. Integration by parts, with the zero trace of $v$,
gives
\[
 \ell(v)-a(w,v)=\int_P(y-\nabla w)\cdot\nabla v,
 \qquad \energy{u-w}\leq\|y-\nabla w\|_{L^2(P)}.
\]
This is a guaranteed estimate: it is valid for every admissible $y$, not
only the minimizer of a flux-fitting problem. It is a functional majorant
in the sense discussed by Repin and by the guaranteed finite element
literature \cite{Repin2008,LiuOishi2013,LiLiu2018}.

A convenient exact construction is
\begin{equation}
 y=-x/2+\curl\psi_h,
 \qquad \curl\psi_h=(\partial_y\psi_h,-\partial_x\psi_h),
 \label{l2:flux0}
\end{equation}
where $\psi_h$ is continuous and piecewise affine. In distributions,
$\operatorname{div}\curl\psi_h=0$. Across an element edge the normal
component of the curl is a tangential derivative of the common trace of
$\psi_h$, so it agrees from both sides. Thus this is a global
$H(\operatorname{div})$ field with precisely the required divergence.

The numerical code chooses $\psi_h$ by minimizing a quadratic mismatch.
Its computed nodal values are simply admissible flux data. An inaccurate
fit widens the bound; it does not destroy equilibration. The first part
$-x/2$ is essential: a curl alone would have the wrong divergence.

\begin{exercise}
What happens if the flux is not exactly equilibrated?
\end{exercise}
\noindent\emph{Answer.} The still-valid estimate is
$\energy{u-w}\leq\|y-\nabla w\|+C_F\|1+\operatorname{div}y\|$,
where $\|v\|\leq C_F\energy v$. Exact equilibration removes the second term.

\section{Differentiate the weak problem, not the approximate output}
\label{l2:pullback}
For two vertex directions, use
$x\mapsto x+s\theta_q(x)+t\theta_r(x)$. The pulled-back coefficient matrix
and volume density are
\[
 \mathbb A(s,t)=\det(I+sD\theta_q+tD\theta_r)
 (I+sD\theta_q+tD\theta_r)^{-1}
 (I+sD\theta_q+tD\theta_r)^{-T},
\]
\[
 a_{s,t}(v,w)=\int_P\nabla v^T\mathbb A(s,t)\nabla w,
 \qquad
 \ell_{s,t}(v)=\int_P\det(I+sD\theta_q+tD\theta_r)v.
\]
Subscripts $q,r,qr$ on forms mean derivatives at $s=t=0$.
In particular,
\begin{align*}
 \mathbb A_q&=(\operatorname{div}\theta_q)I-D\theta_q-D\theta_q^T,\\
 \ell_q(v)&=\int_P(\operatorname{div}\theta_q)v,\\
 \ell_{qr}(v)&=\int_P\big[(\operatorname{div}\theta_q)
 (\operatorname{div}\theta_r)-\tr(D\theta_qD\theta_r)\big]v.
\end{align*}
Differentiating $a_{s,t}(u_{s,t},v)=\ell_{s,t}(v)$ gives
\begin{align}
 a(u_q,v)&=\ell_q(v)-a_q(u,v),\notag\\
 a(u_{qr},v)&=\ell_{qr}(v)-a_q(u_r,v)-a_r(u_q,v)-a_{qr}(u,v).
 \label{l2:exactderivatives}
\end{align}
All derivatives here are material derivatives on the fixed reference fan.
They exist in $H^1_0$ by differentiability of the uniformly coercive
solution operator; that fact alone does not give broken $H^2$ regularity.

Replace the test space by $V_h$ and the state by $u_h$. The exact discrete
first derivatives solve the resulting equations, and $z_h=u_{h,qr}$
solves the fourth discrete equation. Define the continuous lifting by
\begin{equation}
 a(Z_{qr}^h,v)=\ell_{qr}(v)-a_q(u_{h,r},v)
                    -a_r(u_{h,q},v)-a_{qr}(u_h,v).
 \label{l2:lifting}
\end{equation}
Its Galerkin solution is $z_h$. The superscript $h$ on $Z_{qr}^h$ refers
to its data; $Z_{qr}^h$ itself is not a finite element function.

The following form bounds are finite maxima of matrix operator norms:
\[
 |a_q(v,w)|\leq C_q\energy v\energy w,\quad
 |a_r(v,w)|\leq C_r\energy v\energy w,\quad
 |a_{qr}(v,w)|\leq C_{qr}\energy v\energy w.
\]
They transfer errors between the successive equations. No smallness of
$C_q,C_r,C_{qr}$ is assumed.

\section{Why the second variation admits products of errors}
\label{l2:identity}
Set $e=u-u_h$, $e_q=u_q-u_{h,q}$, and $e_r=u_r-u_{h,r}$.
Galerkin orthogonality gives the exact identity
\begin{equation}
 J-J_h=\tfrac12a(e,e).
 \label{l2:energyidentity}
\end{equation}
It holds on nearby polygons after pullback. Differentiating twice means
differentiating the form as well as both arguments:
\[
 (J-J_h)_{qr}=\tfrac12a_{qr}(e,e)+a_q(e_r,e)+a_r(e_q,e)
                      +a(e_q,e_r)+a(e_{qr},e).
\]
Since $u_{h,qr}\in V_h$, orthogonality gives
$a(e_{qr},e)=a(u_{qr},e)$. Insert the second equation in
\eqref{l2:exactderivatives}, tested with $e$, and combine the terms.
The result is
\begin{equation}
 (J-J_h)_{qr}=a(e_q,e_r)-\tfrac12a_{qr}(e,e)
                         +a(Z_{qr}^h-z_h,e).
 \label{l2:productidentity}
\end{equation}
The exact second derivative has disappeared. If
\[
 \energy e\leq\delta_0,\quad \energy{e_q}\leq\delta_q,
 \quad\energy{e_r}\leq\delta_r,\quad
 \energy{Z_{qr}^h-z_h}\leq\eta_{qr},
\]
then
\begin{equation}
 |(J-J_h)_{qr}|\leq
 \delta_q\delta_r+\tfrac12C_{qr}\delta_0^2+\delta_0\eta_{qr}.
 \label{l2:productbound}
\end{equation}
This is manuscript equation~\meq{eq:second-residual-bound}.
It is valid on each fixed mesh without an asymptotic regularity assumption.

\subsection{Could we use only material derivatives?}
Yes. The weak equation already gives
$a(u_q,u_r)=\ell_q(u_r)-a_q(u,u_r)$. It uses no second derivative.
But replacing $u$ by $u_h$ changes the load, so it does not automatically
remove the terms linear in the state error. Direct propagation of
$O(h)$ energy bounds through the Hessian generally leaves an $O(h)$ bound.
The Bogosel--Bucur analysis already exploits a quadratic interpolation
product for material derivatives; the leading first-order contribution
comes from other data errors, not a failure of Galerkin orthogonality
\cite[Section 3.1]{BogoselBucurValidated2024}.

In \eqref{l2:productbound}, all four errors of order $h$ instead produce
an order-$h^2$ bound. This improvement is due to the cancellation and
lifting. Equilibrated fluxes are a way to enclose the energy errors;
they are not, by themselves, a theorem of quadratic convergence.

\section{The computed linear systems also need error bounds}
\label{l2:algebraic}
Let $Kx_*=f$ be the exact interior-node system, with $K$ symmetric positive
definite. A computed vector $x$ has residual $f-Kx$. If
$\lambda_{\min}(K)\geq\underline\alpha>0$, then
\[
 \|x-x_*\|_K^2=(f-Kx)^TK^{-1}(f-Kx)
 \leq \frac{\|f-Kx\|_2^2}{\underline\alpha}.
\]
The Euclidean norm of the residual is not the energy norm of the error;
the factor $1/\sqrt{\underline\alpha}$ supplies that conversion.

For the circumradius-one polygon, domain monotonicity and the square give
$\lambda_1(P)\geq\pi^2/2$. On one triangle,
\[
 M_T=\frac{|T|}{12}
 \begin{pmatrix}2&1&1\\1&2&1\\1&1&2\end{pmatrix}
 \succeq\frac{|T|}{12}I.
\]
After assembly and restriction to interior nodes,
\begin{equation}
 \underline\alpha=
 \frac{\pi^2}{24}\min_{i\text{ interior}}\sum_{T\ni i}|T|
 \quad\hbox{is a valid lower bound for }\lambda_{\min}(K).
 \label{l2:coercivity}
\end{equation}
This is explicit and possibly conservative. It contains no numerically
estimated inverse and no uncertified eigenvalue computation.

Define the four residual functionals of the computed candidates:
\begin{align*}
 \rho_0(v)&=\ell(v)-a(\widetilde u,v),\\
 \rho_q(v)&=\ell_q(v)-a_q(\widetilde u,v)-a(\widetilde u_q,v),\\
 \rho_r(v)&=\ell_r(v)-a_r(\widetilde u,v)-a(\widetilde u_r,v),\\
 \rho_z(v)&=\ell_{qr}(v)-a_q(\widetilde u_r,v)
 -a_r(\widetilde u_q,v)-a_{qr}(\widetilde u,v)-a(\widetilde z,v).
\end{align*}
Their vectors are obtained by evaluating them on each interior basis
function $\varphi_i$: $(r_\rho)_i=\rho(\varphi_i)$. For example,
$r_{\rho_0}=f-K\widetilde x$.

Because later loads contain earlier unknowns, the transfer is
\begin{align*}
 \epsilon_0&=\|r_{\rho_0}\|_2/\sqrt{\underline\alpha},\\
 \epsilon_q&=\|r_{\rho_q}\|_2/\sqrt{\underline\alpha}+C_q\epsilon_0,\\
 \epsilon_r&=\|r_{\rho_r}\|_2/\sqrt{\underline\alpha}+C_r\epsilon_0.
\end{align*}
These bound the differences between exact discrete solutions and their
candidates. Boundary coefficients are set to exactly zero before any
residual is assembled. A large penalty imposed by a floating-point
solver is not a substitute for this Dirichlet restriction.

\subsection{Correct the centre before increasing the radius}
The energy centre
$\widetilde J_{\rm corr}=\ell(\widetilde u)-a(\widetilde u,\widetilde u)/2$
satisfies $J_h-\widetilde J_{\rm corr}=\energy{u_h-\widetilde u}^2/2$.
For the mixed derivative use
\begin{align*}
 \widetilde C_{\rm corr}={}&\ell_{qr}(\widetilde u)
 +\ell_q(\widetilde u_r)-\tfrac12a_{qr}(\widetilde u,\widetilde u)
 -a_q(\widetilde u,\widetilde u_r)\\
 &+\rho_r(\widetilde u_q)+\rho_0(\widetilde z).
\end{align*}
Expanding about the exact discrete fields and substituting the residual
identities gives
\begin{align*}
 J_{h,qr}-\widetilde C_{\rm corr}
 ={}&\rho_z(u_h-\widetilde u)
 -\tfrac12a_{qr}(u_h-\widetilde u,u_h-\widetilde u)\\
 &+a(u_{h,q}-\widetilde u_q,u_{h,r}-\widetilde u_r).
\end{align*}
Every first-order material-solve term has cancelled. Consequently
\[
 |J_{h,qr}-\widetilde C_{\rm corr}|\leq E_{\rm alg}
 :=\frac{\|r_{\rho_z}\|_2}{\sqrt{\underline\alpha}}\epsilon_0
   +\tfrac12C_{qr}\epsilon_0^2+\epsilon_q\epsilon_r.
\]
The lifting solve error need not be given another separate name: its
residual appears directly in this expression.

\section{Four fluxes and the final continuous enclosure}
\label{l2:assembly}
Choose four continuous affine potentials. The state flux is
\eqref{l2:flux0}; the first-derivative flux is
$-\theta_q+\curl\psi_{q,h}$, with divergence $-\operatorname{div}\theta_q$.
For the second lifting use
\[
 -(\operatorname{div}\theta_q)\theta_r+D\theta_q\theta_r
                                  +\curl\psi_{z,h}.
\]
Indeed,
\[
 \operatorname{div}\big((\operatorname{div}\theta_q)\theta_r
                 -D\theta_q\theta_r\big)
 =(\operatorname{div}\theta_q)(\operatorname{div}\theta_r)
                       -\tr(D\theta_qD\theta_r).
\]
Its normal trace is continuous across fan edges: the gradient jump of a
continuous affine field is a vector tensor the edge normal, and the two
normal jump contributions cancel. This is a distributional identity,
not only a trianglewise identity.

The four mismatch norms are
\begin{align*}
 \Phi_0&=\|-x/2+\curl\psi_{0,h}-\nabla\widetilde u\|,\\
 \Phi_q&=\|-\theta_q+\curl\psi_{q,h}-\nabla\widetilde u_q
                                     -\mathbb A_q\nabla\widetilde u\|,\\
 \Phi_r&=\|-\theta_r+\curl\psi_{r,h}-\nabla\widetilde u_r
                                     -\mathbb A_r\nabla\widetilde u\|,\\
 \Phi_z&=\|-(\operatorname{div}\theta_q)\theta_r+D\theta_q\theta_r
                 +\curl\psi_{z,h}-\nabla\widetilde z\\
 &\hspace{23mm}-\mathbb A_q\nabla\widetilde u_r
                 -\mathbb A_r\nabla\widetilde u_q
                 -\mathbb A_{qr}\nabla\widetilde u\|.
\end{align*}
All unmarked norms in these formulas are $L^2(P)$ norms of vector fields.
Residual bounds, load-error transfer, and best approximation give
\begin{align*}
 \widehat\delta_0&=\Phi_0,\\
 \widehat\delta_q&=\Phi_q+C_q\Phi_0+\epsilon_q,
 &\widehat\delta_r&=\Phi_r+C_r\Phi_0+\epsilon_r,\\
 \widehat\eta_{qr}&=\Phi_z+C_q\epsilon_r+C_r\epsilon_q+C_{qr}\epsilon_0.
\end{align*}
For example, $\Phi_0$ bounds the error to $\widetilde u$; best
approximation then also bounds the error to $u_h$. In the lifting bound,
$\widetilde z\in V_h$ can replace $z_h$ before best approximation is used.
This explains why an extra lifting algebraic error is not added there.

For a nonzero Fourier pairing, $A_q=A_r=0$. Define
\[
 B_J=\widehat\delta_q\widehat\delta_r+\tfrac12C_{qr}\widehat\delta_0^2
                           +\widehat\delta_0\widehat\eta_{qr}.
\]
The entry centre and radius are
\begin{align}
 \widetilde F_{qr}&=\frac{\widetilde C_{\rm corr}}{A^2}
                     -\frac{2\widetilde J_{\rm corr}A_{qr}}{A^3},\notag\\
 R_F&=\frac{B_J+E_{\rm alg}}{A^2}
       +\frac{|A_{qr}|}{A^3}(\widehat\delta_0^2+\epsilon_0^2),\notag\\
 |D^2\F[q,r]-\widetilde F_{qr}|&\leq R_F.
 \label{l2:finalradius}
\end{align}
Every term has a job: $B_J$ is the continuous discretization error,
$E_{\rm alg}$ is the error in the discrete centre, and the last term
accounts for the area-normalization coefficient multiplying $J$.
If the evaluated centre is an interval, its arithmetic radius is retained.

\section{Regularity attempt: integrate the vertex-zero load by parts}
\label{l2:rayproof}
This section develops the proposed Bogosel--Bucur route. We will prove
broken regularity of first and second material derivatives \emph{at the
regular polygon}, using the standard polygonal mixed-boundary regularity
theorem stated below. A separate final paragraph explains why this does
not yet prove the manuscript's stronger uniform-neighborhood formulation.

For $c=1,2$, put the weak equation for $U_0^c$ in the form
\begin{align}
 a(U_0^c,v)&=\int_P G_{0,c}(u)\cdot\nabla v
                              +\int_P(\partial_c\phi_0)v,\notag\\
 G_{0,c}(u)&=-(\partial_c\phi_0)\nabla u
            +(\partial_cu)\nabla\phi_0
            +(\nabla\phi_0\cdot\nabla u)e_c.
 \label{l2:G}
\end{align}
On a sector the hat gradient is constant. Since $\Delta u=-1$,
\[
 \operatorname{div}G_{0,c}(u)
 =\partial_c\phi_0+2\nabla\phi_0\cdot\nabla\partial_cu.
\]
The volume load after integration by parts is therefore
\begin{equation}
 -2\nabla\phi_0\cdot\nabla\partial_cu\in L^2(P).
 \label{l2:regularload}
\end{equation}
It is supported on $T_{-1}\cup T_0$. The term
$\partial_c\phi_0$ has cancelled; keeping it would give an incorrect
strong formulation.

Let $S_j=\{r a_j:0<r<1\}$ and let
$\nu_j=(-\sin(j\vartheta),\cos(j\vartheta))$ point from $T_{j-1}$ into
$T_j$. The trace convention is $[w]_j=w|_{T_j}-w|_{T_{j-1}}$.
The contribution of the two boundary integrals is
$-[G_{0,c}(u)]_j\cdot\nu_j$.
Continuity of $\phi_0$ implies $[\nabla\phi_0]_j$ is normal to $S_j$.
Because $u\in H^2(P)$, its gradient has a common trace there. Substituting
into \eqref{l2:G} gives the especially simple ray density
\begin{equation}
 -[G_{0,c}(u)]_j\cdot\nu_j
     =-([\nabla\phi_0]_j\cdot\nu_j)\,\partial_cu|_{S_j}.
 \label{l2:raydensity}
\end{equation}
Only $S_{-1},S_0,S_1$ occur. Their coefficients are
\[
 [\nabla\phi_0]_0\cdot\nu_0=-2\cot\vartheta,
 \qquad
 [\nabla\phi_0]_{\pm1}\cdot\nu_{\pm1}=\frac1{\sin\vartheta}.
\]
Thus the complete distributional equation is
\begin{equation}
 -\Delta U_0^c=-2\nabla\phi_0\cdot\nabla\partial_cu
 +2\cot\vartheta\,(\partial_cu)|_{S_0}\delta_{S_0}
 -\frac{(\partial_cu)|_{S_1}\delta_{S_1}
        +(\partial_cu)|_{S_{-1}}\delta_{S_{-1}}}{\sin\vartheta}.
 \label{l2:decomposition}
\end{equation}
Here $g\delta_S$ means the functional $v\mapsto\int_Sgv\,ds$.
No point masses occur: we integrated bounded piecewise coefficients times
$H^1$ gradient fields once, and the singular terms are edge measures.

\subsection{Why the ray densities have the right endpoint behavior}
For fixed $n$, convex-polygon Dirichlet regularity gives
\[
 u\in H^{2+\sigma}(P),\qquad
 0<\sigma<\min\left\{\frac12,\frac{\pi}{\pi-2\pi/n}-1\right\},
\]
after decreasing $\sigma$ if necessary. Possible resonant logarithmic
corner terms are harmless below this strict exponent threshold.
In particular $\nabla u$ is continuous up to each corner. Its tangential
derivatives on the two adjacent Dirichlet edges are zero; the tangents
are independent, so $\nabla u(a_j)=0$. At the centre, rotational symmetry
implies $\nabla u(0)=0$.

On $S_j$, reflection gives zero normal derivative, hence
\[
 \nabla u(r a_j)=a_j\,\partial_ru(r,0).
\]
Thus the densities in \eqref{l2:decomposition} are precisely rotated
multiples of the same radial derivative. The trace theorem gives
$H^{1/2+\sigma}$ regularity on each ray, and the values at $r=0,1$ vanish.
Their extensions by zero across the centre belong to the same fractional
space for sufficiently small $\sigma>0$. Endpoint values here are justified
by the extra positive regularity; bare $H^{1/2}$ membership would not by
itself justify a pointwise endpoint argument.

\subsection{The mixed-boundary ingredient, with its hypotheses visible}
\begin{lemma}[A ray load with vanishing endpoints]
\label{l2:raylemma}
Fix a regular polygon. Let $g\in H^{1/2+\sigma}(0,1)$ with $g(0)=g(1)=0$
and $\sigma>0$ sufficiently small. The solution $W\in H^1_0(P)$ of
$a(W,v)=\int_{S_0}gv$ belongs to $H^{2+\sigma'}$ on each side of the ray
and on every fan sector, for some $\sigma'>0$ that may be smaller than
$\sigma$. The same assertion holds for any rotated ray.
\end{lemma}
\begin{proof}
The line load is continuous on $H^1_0(P)$ by the trace theorem, so
Lax--Milgram gives $W$. Its data and domain are invariant under reflection
across the line through $S_0$, hence $W$ is even. Restrict to the upper
half-polygon. The boundary conditions are Dirichlet zero on the polygon
boundary and outward Neumann data $g/2$ on the positive part of the cut,
zero on its negative part. The factor $1/2$ follows by doubling the weak
integral over the half-polygon.

Apply the polygonal Dirichlet--Neumann shift theorem of Grisvard
\cite[Theorem 4.4.3.7 and Corollary 4.4.3.8]{Grisvard1985}.
Its use here is the same as the
ray decomposition in \cite[Section 3]{BogoselBucurValidated2024}; the local
reviewed PDF has Lemma 3.1 and Remark 3.2. The hypotheses can be checked
geometrically. At a vertex on the cut the mixed angle is
$(\pi-\vartheta)/2<\pi/2$, whose first mixed corner exponent is
$\pi/(\pi-\vartheta)>1$. Other Dirichlet angles are less than $\pi$.
When $n$ is odd, the negative cut meets the opposite edge at a right angle;
the Neumann data are identically zero near that point, and even reflection
reduces the local problem to a straight Dirichlet boundary. There is no
nonpolynomial corner singularity of exponent at most one there.
The Neumann datum extends through the centre in $H^{1/2+\sigma}$,
and its zero value at the positive endpoint supplies the compatibility
with the zero Dirichlet data. These facts give a positive shift above
$H^2$, after shrinking the exponent. Reflection back gives the assertion.
Away from the loaded positive ray, the equation is locally harmonic across
the artificial negative cut, so no extra interface is introduced.
\end{proof}

\begin{proposition}[First material derivatives at the regular polygon]
\label{l2:firstregular}
For every fixed $n\geq3$, the first vertex material derivatives are in
$H^{2+\sigma}$ on each coarse fan triangle for some $\sigma>0$ depending
on $n$. Their gradients have a common value at the centre and vanish at
the polygon vertices.
\end{proposition}
\begin{proof}
Split \eqref{l2:decomposition} into its volume solution and three ray
solutions. The volume load is piecewise $H^\sigma$, hence globally
$H^\sigma$ after choosing $\sigma<1/2$; multiplication by a sector
indicator preserves $H^\sigma$ in this range. The convex Dirichlet
polygonal shift theorem gives $H^{2+\sigma}$ for its solution, with a
possibly smaller positive exponent. Lemma~\ref{l2:raylemma} handles the
three ray solutions. Rotation and linear combination handle all vertices
and all directions $q$.

The one-sided gradients are continuous to sector endpoints by Sobolev
embedding. Across a ray, the tangential gradient agrees, because the
function has one common trace. The normal-gradient jump equals the line
density, up to the chosen orientation. That density vanishes at both
endpoints, so the gradients agree at the centre and at each outer endpoint.
At a polygon vertex the common gradient is orthogonal to both adjacent
boundary tangents, hence is zero. At the centre the common gradient need
not be zero; symmetry of the state does not imply the same symmetry for
a single-vertex material derivative.
\end{proof}

\section{The second derivatives: a necessary centre cancellation}
\label{l2:secondregular}
The first-derivative proof is not automatically a proof for $u_{qr}$.
Nevertheless it supplies the extra regularity needed to continue. On each
sector write its weak equation as
\[
 a(u_{qr},v)=\int_P d_{qr}v+\int_P G_{qr}\cdot\nabla v,
\]
where, for this regularity calculation only,
\begin{align*}
 d_{qr}&=(\operatorname{div}\theta_q)(\operatorname{div}\theta_r)
                               -\tr(D\theta_qD\theta_r),\\
 G_{qr}&=-\mathbb A_q\nabla u_r-\mathbb A_r\nabla u_q
                                  -\mathbb A_{qr}\nabla u.
\end{align*}
These two names display the volume and flux parts of one distribution;
they do not introduce any new PDE solves. Sectorwise integration by parts
produces $d_{qr}-\operatorname{div}G_{qr}\in L^2$ and ray densities
$-[G_{qr}]_j\cdot\nu_j\in H^{1/2+\sigma}(S_j)$.
All outer endpoint values vanish by Proposition~\ref{l2:firstregular}.
The centre values may not vanish because $\nabla u_q(0)$ and
$\nabla u_r(0)$ may be nonzero.

This matters. A constant load on a ray ending at an interior point can
produce an $r\log r$ corner term; its second derivatives are not square
integrable near that endpoint. Splitting a compatible collection of rays
into unrelated single-ray problems can therefore destroy the regularity
one is trying to prove.

\subsection{An explicit local lifting removes the centre values}
Choose a smooth cutoff $\chi$ supported in an interior disk around the
centre, equal to one in a smaller disk. Subtract from $u_{qr}$ the known
piecewise smooth function
\begin{equation}
 \chi(x)\bigl(\theta_q(x)\cdot\nabla u_r(0)
                  +\theta_r(x)\cdot\nabla u_q(0)\bigr).
 \label{l2:centerlifting}
\end{equation}
It belongs to $H^1_0(P)$ and is $H^2$ on each fan sector. Its purpose is
local regularity analysis; it is not the numerical second lifting $Z_{qr}^h$.
The centre gradients are well defined by the first-derivative proposition.

For any constant vector $b$, the following identity holds in distributions:
\begin{equation}
 \operatorname{div}(-\mathbb A_q b)
     =\Delta(\theta_q\cdot b).
 \label{l2:centeridentity}
\end{equation}
Indeed,
$-\mathbb A_qb=D\theta_qb+D\theta_q^Tb-(\operatorname{div}\theta_q)b$.
The divergences of the first and last terms cancel by commutation of
weak derivatives; the middle term is $\nabla(\theta_q\cdot b)$.
Near the centre, where $\chi=1$, identity~\eqref{l2:centeridentity} shows
that \eqref{l2:centerlifting} has exactly the ray-jump contributions of
$-\mathbb A_q\nabla u_r(0)-\mathbb A_r\nabla u_q(0)$.
Also $\nabla u(0)=0$, so the $\mathbb A_{qr}$ term contributes no centre
value. Subtracting the lifting therefore makes every remaining ray
density vanish at the centre.

Differentiating the cutoff introduces only piecewise smooth volume terms
and ray terms supported away from the centre and the outer vertices.
The resulting volume source is still $L^2$; each ray density is still
$H^{1/2+\sigma}$ and now vanishes at \emph{both} endpoints. The regular
volume problem and the isolated ray problems can now be treated separately.

\begin{proposition}[Second material derivatives at the regular polygon]
\label{l2:secondproposition}
Under the polygonal shift theorem used in Lemma~\ref{l2:raylemma}, the
exact second material derivative $u_{qr}$ belongs to $H^2(T_j)$ for every
fan triangle, every fixed regular polygon, and every pair of vertex
directions $q,r$.
\end{proposition}
\begin{proof}
After subtracting \eqref{l2:centerlifting}, solve the $L^2$ volume problem
with zero Dirichlet data. Its solution is globally $H^2$ by convexity.
For each remaining ray load apply Lemma~\ref{l2:raylemma}. The sum of
these solutions and the subtracted lifting satisfies the weak equation
for $u_{qr}$, and uniqueness identifies it with $u_{qr}$. Every summand
is $H^2$ on each fan triangle. There are finitely many rays and finitely
many vertex coordinates, so the argument applies to all mixed directions.
\end{proof}

\subsection{What part of the conjecture has this proved?}
The argument proves the qualitative broken regularity of the state,
first derivatives, and second derivatives \emph{at $P_n$}, for fixed $n$.
It also gives uniformity over bounded vertex directions there, by finite
dimensionality and bilinearity. Constants may depend on $n$; the corner
exponent tends to one as $n$ grows.

The manuscript's Conjecture~\meq{conj:h2} asks for uniform broken
regularity for the parameterized family in a neighborhood. That stronger
statement does not follow just by rotating and reflecting the regular
polygon: the perturbed polygons no longer have its reflection symmetries.
A proof still needs stability of the relevant transmission estimates
under these geometric perturbations, or a parameter-dependent corner
expansion with uniform control. These notes do not claim that step is
proved, and do not change the manuscript's conjecture or its theorem.

For the Hessian rate \emph{at the regular polygon}, the regularity just
proved is sufficient. Broken interpolation and the differentiated
Galerkin equations give $\delta_0,\delta_q,\delta_r=O(h)$. Moreover,
subtracting the exact second-derivative equation from \eqref{l2:lifting}
gives
\[
 \energy{Z_{qr}^h-u_{qr}}
 \leq C_q\delta_r+C_r\delta_q+C_{qr}\delta_0=O(h).
\]
Approximate the fixed function $u_{qr}$ by its broken interpolant and use
best approximation for $Z_{qr}^h$ to obtain $\eta_{qr}=O(h)$.
Equation~\eqref{l2:productbound} then gives
$D^2J(P_n)-D^2J_h(P_n)=O(h^2)$.
The energy identity and its first derivative also bound $J-J_h$ and
$D(J-J_h)$ by $O(h^2)$, so differentiating $A^{-2}$ gives the same rate
for the normalized Hessian. This is a statement about exact Galerkin
quantities, not yet about the computed radius in \eqref{l2:finalradius}.

\subsection*{Exercises and proof checkpoints}
\begin{exercise}
For $c=2$, simplify the three ray terms in \eqref{l2:decomposition} using
$\nabla u(r a_j)=a_j\partial_ru(r,0)$.
\end{exercise}
\noindent\emph{Answer.} The $S_0$ term is zero. The terms on $S_1,S_{-1}$
are respectively $-\partial_ru\,\delta_{S_1}$ and
$+\partial_ru\,\delta_{S_{-1}}$, with the traces understood on their rays.
For $c=1$ the weights are $2\cot\vartheta,-\cot\vartheta,-\cot\vartheta$.

\begin{exercise}
Why is pointwise zero at an endpoint not an argument for a general
$H^{1/2}$ function?
\end{exercise}
\noindent\emph{Answer.} Endpoint evaluation is not a continuous trace
at the borderline exponent. The positive extra regularity above, or an
explicit weighted endpoint condition, must be used.

\begin{exercise}
Which extra facts are needed to conclude that the \emph{computed}
$R_F$ is $O(h^2)$?
\end{exercise}
\noindent\emph{Answer.} The four chosen flux mismatches and the transferred
algebraic bounds must have compatible orders, so that $B_J$, $E_{\rm alg}$,
and the normalization remainder are $O(h^2)$. A fixed solver tolerance or
fixed precision can eventually prevent that asymptotic behavior even when
the exact Galerkin Hessian converges quadratically.
